\section{Билет 51. Вероятность. Среднее значение. Функция распределения вероятности. Условие нормировки. Вычисление средних величин.}

\begin{definition}
\textbf{Вероятностью} $P_i$ того, что случайная величина $x$ примет значение $x_i$, называется отношение числа благоприятных исходов $N_i$ к общему числу испытаний $N$:
\begin{equation}
    P_i = \frac{N_i}{N}. \label{eq:prob_def_51}
\end{equation}
Вероятность --- безразмерная величина, принимающая значения от 0 до 1. Сумма вероятностей всех возможных значений случайной величины равна единице:
\begin{equation}
    \sum_i P_i = 1. \label{eq:prob_sum_51}
\end{equation}
\end{definition}

\begin{definition}
\textbf{Средним значением} $\langle x \rangle$ случайной величины $x$ называется взвешенная сумма всех возможных значений, где весами служат соответствующие вероятности:
\begin{equation}
    \langle x \rangle = \sum_i x_i P_i. \label{eq:mean_discrete_51}
\end{equation}
Для непрерывной случайной величины, принимающей значения в интервале $(a, b)$, среднее значение вычисляется через функцию распределения:
\begin{equation}
    \langle x \rangle = \int_a^b x \, f(x) \, dx, \label{eq:mean_continuous_51}
\end{equation}
где $f(x)$ --- плотность вероятности.
\end{definition}

\begin{definition}
\textbf{Функцией распределения вероятности} (плотностью вероятности) $f(x)$ называется величина, определяемая как предел отношения вероятности попадания случайной величины в интервал $(x, x+dx)$ к ширине этого интервала:
\begin{equation}
    f(x) = \frac{dP}{dx}, \quad \text{или} \quad dP = f(x) \, dx. \label{eq:distribution_func_51}
\end{equation}
Физический смысл $f(x)$: вероятность того, что величина $x$ окажется в единичном интервале вблизи значения $x$.
\end{definition}

\begin{theorem}[Условие нормировки]
Поскольку случайная величина обязательно принимает какое-либо значение из всего множества возможных, полная вероятность равна единице. Это даёт \textbf{условие нормировки} для функции распределения:
\begin{equation}
    \int_{-\infty}^{+\infty} f(x) \, dx = 1. \label{eq:normalization_51}
\end{equation}
Для дискретного распределения условие нормировки имеет вид $\sum_i P_i = 1$.
\end{theorem}

\begin{property}[Вычисление средних величин]
Для любой функции $\varphi(x)$ от случайной величины $x$ среднее значение вычисляется по формуле:
\begin{equation}
    \langle \varphi(x) \rangle = \int_{-\infty}^{+\infty} \varphi(x) \, f(x) \, dx. \label{eq:mean_general_51}
\end{equation}
Важные частные случаи:
\begin{itemize}
    \item Среднее значение самой величины: $\langle x \rangle = \int x f(x) dx$;
    \item Средний квадрат: $\langle x^2 \rangle = \int x^2 f(x) dx$;
    \item Дисперсия: $D[x] = \langle x^2 \rangle - \langle x \rangle^2$.
\end{itemize}
\end{property}

\begin{remark}
\textbf{Теорема об умножении вероятностей.} Если две случайные величины $x$ и $y$ статистически независимы, то вероятность их одновременного появления равна произведению вероятностей:
\begin{equation}
    P(x_i, y_j) = P(x_i) \cdot P(y_j). \label{eq:independence_51}
\end{equation}
Это свойство широко используется в статистической физике, например, при выводе распределения Максвелла, где проекции скорости $v_x$, $v_y$, $v_z$ считаются независимыми.
\end{remark}
